\section{Introduction}\label{sec:intro}

Unmanned aerial vehicles (UAVs) are widely used for traffic monitoring,
surveillance, and search, and many of these tasks need multi-object tracking
(MOT): every object must be detected in every frame and keep the same
identity over time \cite{wu2022uavsurvey,zhu2022visdrone}. Aerial video makes
this hard. Objects are small, their number changes quickly, the camera moves,
and illumination and blur change within one video
\cite{zhu2022visdrone,du2018uavdt,cheng2023smallobject}. Many of these
applications also need results with low delay, so the processing cost matters
as much as the accuracy.

Most MOT systems follow the tracking-by-detection design. A detector finds the
objects in each frame, and a tracker links the detections over time
\cite{bewley2016sort,zhang2022bytetrack}. Much research improves the tracker,
for example its motion model or its association rules, and much research
improves the detector itself. Far less attention is given to how the detector
is run. In practice the detector uses one confidence threshold, one
non-maximum-suppression (NMS) threshold, and one input resolution for the
whole video. These three values form the detector operating point, and they
decide which boxes reach the tracker.

A fixed operating point does not fit a scene that keeps changing. A setting
that works in a sparse, clear scene can keep too much clutter in a dense
scene. A setting that works in a dense scene can lose small or weak objects in
another part of the same video. A large input helps small objects, but it
costs the same time on easy and hard frames. The choice therefore moves false
positives, missed objects, identity switches, and runtime together.

Existing adaptive methods only partly address this problem. Adaptive inference
methods change the detector's input scale or configuration, but they usually
need an extra learned model and are designed for detection, not tracking
\cite{chin2019adascale,xu2020approxdet}. Adaptive tracking methods change
thresholds inside one specific tracker \cite{ma2023adaptivebytetrack,shim2024adaptrack}.
What is missing is a simple, training-free way to adapt the detector
operating point online for MOT, without changing the detector or the tracker.

This paper presents AC-MOT, a training-free control layer that adapts the
detector operating point online for aerial MOT. It sits between the detector
and the tracker, reads inexpensive scene cues, and works with existing
detectors and trackers without retraining or modifying them. The method is
described in Sect.~\ref{sec:method}, and Sect.~\ref{sec:development} explains
how it was developed with optimization and tested step by step.

The main contributions are:
\begin{enumerate}
\item A training-free control layer that treats the detector operating point
      as part of the MOT system and adapts it online (Sect.~\ref{sec:method}).
\item An optimization-guided development path, in which each step is compared
      with the baseline it was built on (Sect.~\ref{sec:development}).
\item An evaluation on VisDrone2019-MOT with two trackers, ByteTrack and the
      recent OATrack, using eight MOT metrics and paired bootstrap intervals,
      together with a transfer to UAVDT without retuning
      (Sects.~\ref{sec:setup} and~\ref{sec:results}).
\item A latency measurement showing that the control layer adds only a few
      milliseconds per frame on a Tesla T4 GPU (Sect.~\ref{sec:discussion}).
\end{enumerate}

The rest of the paper is organized as follows. Section~\ref{sec:related}
reviews related work. Section~\ref{sec:method} describes AC-MOT, and
Section~\ref{sec:development} explains how it was developed.
Section~\ref{sec:setup} gives the experimental setup,
Section~\ref{sec:results} the results, and Section~\ref{sec:discussion} a
discussion. Section~\ref{sec:conclusion} concludes the paper.
